\section{Generalized proof tools}
\label{chap:generalized-tools}

This section develops the generalized proof rules that consume the comparison
foundations established earlier.

\subsection{Generalized Leibniz and Mahowald rules}

% SOURCE: MainPaper/main.tex:1606-1652, Theorem thm:e73f481e.
\begin{theorem}[Generalized Leibniz rule]
  \label{thm:generalized-leibniz}
  \uses{prop:delta-ess-formula-infinite,thm:delta-classical-equivalence,prop:classical-crossing-iff-delta-crossing,def:page-extension,prop:page-crossing-iff-synthetic-crossing,cor:fess-square-naturality}
  \lean{KIP126.Kervaire.Route.Tools.GeneralizedLeibnizLaw}
  \notready
  Let $f:X\to Y$, $2\le n\le r$,
  $e(f)\le m\le n-2+e(f)$, and $l\ge e(f)$.  Suppose
  \[
  \begin{aligned}
    &x\in Z_{r-1}^{s,t}(X),&&
    y\in Z_{r-1-m+e(f)}^{s+m,t+m}(Y),\\
    &x_\infty\in Z_\infty^{s+r,t+r-1}(X),&&
    y_\infty\in Z_\infty^{s+r+l,t+r+l-1}(Y),
  \end{aligned}
  \]
  and assume:
  \begin{enumerate}
    \item $d_r(x)=x_\infty$;
    \item $d_m^{f,E_n}(x)=y$;
    \item $d_l^{f,E_\infty}(x_\infty)=y_\infty$;
    \item either the first differential has no crossing on $E_n$, or the
      second page extension has no crossing;
    \item the third page extension has no crossing.
  \end{enumerate}
  Then $d_{r+l-m}(y)=y_\infty$.
  The linked Lean declaration specifies this law for a fixed actual
  normalized page family; its validity remains to be established.
  It does not assert that arbitrary comparison isomorphisms satisfy the law.
\end{theorem}

\begin{proof}
  Form the square whose horizontal maps are
  $\widehat f_{n-1}$ and $\widehat f$ and whose vertical maps are the two
  connecting maps from reduction modulo $\lambda^{n-1}$.  Translate (1) by the
  $\delta$ formula and (2)--(3) by the definition of page extension.  The two
  crossing comparison theorems turn hypothesis (4) into synchronized-source
  no-crossing and hypothesis (5) into target no-crossing.  Apply ESS square
  propagation.  Finally use the $\delta$--classical equivalence to remove the
  displayed $\lambda$ powers and obtain the asserted classical differential
  coset.
\end{proof}

% SOURCE: MainPaper/main.tex:1664-1692, Example exam:Leibnizyes.
\begin{proposition}[Positive Leibniz regression]
  \label{prop:leibniz-positive-regression}
  \uses{thm:generalized-leibniz,thm:external-adams-one-line,prop:hopf-e2-page-extension,prop:hopf-einfty-page-extension,cor:classical-e2-no-crossing,prop:page-crossing-regression}
  \notready
  The inputs
  $d_2(h_5)=h_0h_4^2$,
  $d_1^{\nu_{\mathrm{Hopf}},E_2}(h_5)=h_2h_5$, and
  $d_2^{\nu_{\mathrm{Hopf}},E_\infty}(h_0h_4^2)=h_0p$
  satisfy the generalized Leibniz hypotheses and give
  \[
    d_3(h_2h_5)=h_0p.
  \]
\end{proposition}

\begin{proof}
  Substitute $(n,r,m,l,e(f))=(2,2,1,2,1)$.  The $E_2$ classical crossing
  condition is automatic, and the two Hopf-map page extensions have the
  no-crossing certificates from the regression.  The theorem gives a
  differential of length $2+2-1=3$ with the claimed target.
\end{proof}

% SOURCE: MainPaper/main.tex:1694-1720, Example exam:Leibnizno.
\begin{proposition}[No-crossing is necessary]
  \label{prop:leibniz-negative-regression}
  \uses{thm:generalized-leibniz,prop:two-page-extension-regression,prop:page-crossing-regression}
  \notready
  For $f=2$ and the computed differential
  $d_2(h_4)=h_0h_3^2$, all generalized Leibniz hypotheses except target
  no-crossing can be met using
  $d_2^{2,E_\infty}(h_0h_3^2)=0$.  The missing hypothesis is false because
  $d_1^{2,E_\infty}(d_0)=h_0d_0$ crosses it.  Omitting that hypothesis would
  assert the false zero differential $d_3(h_0h_4)=0$.
\end{proposition}

\begin{proof}
  Evaluate every hypothesis using the page-extension regression and the typed
  classical computation.  The length-one extension meets the crossing bounds,
  so the generalized theorem is inapplicable.  The known nonzero differential
  on $h_0h_4$ supplies the regression check that the weakened statement is
  false.
\end{proof}

% SOURCE: MainPaper/main.tex:1733-1753, Remark rem:chuaerror; Chua Definition 12.5 and Theorem 12.9.
\begin{proposition}[Counterexample to the maximal-extension rule]
  \label{prop:chua-rule-counterexample}
  \uses{prop:leibniz-negative-regression}
  \notready
  The maximal-$[h_0]$ lift of $h_0h_3^2$ in the stated synthetic $14$-stem
  example maps to zero, while another lift maps to $[h_0d_0]$.  Consequently
  the rule quoted as Chua Theorem 12.9, when read without a no-crossing
  hypothesis, predicts the false conclusion of
  Proposition~\ref{prop:leibniz-negative-regression}.
\end{proposition}

\begin{proof}
  Use the two distinct detected lifts and their products from the supplied
  evidence.  The zero product is maximally $\lambda$-divisible under the cited
  definition.  Substitution into the quoted rule yields the false zero target,
  so these data form a concrete counterexample; no claim about a corrected
  version is made.
\end{proof}

% SOURCE: MainPaper/main.tex:1780-1917, complete hypotheses and degree package for Theorem thm:158d451a.
\begin{definition}[Hypotheses of the generalized Mahowald rule]
  \label{def:mahowald-input-data}
  \uses{prop:normalized-synthetic-triangle,def:page-extension,def:classical-differential-crossing}
  \notready
  The Mahowald hypotheses specify a distinguished triangle
  $X\xrightarrow fY\xrightarrow gZ\xrightarrow h\Sigma X$ with
  $e(f)+e(g)+e(h)=1$, integers $s,t,n,m,l$, and
  \[
    r=n+m+l,\quad n_1=n-e(f)\ge1,\quad
    m_1=m-e(g)\ge0,\quad l_1=l-e(h)\ge0,\quad
    r'=r-m_1=n_1+l_1+1.
  \]
  It contains classes
  \[
  \begin{aligned}
    x&\in Z_{n_1}^{s+l,t+l-1}(X),&
    y&\in Z_{m_1+1}^{s+n+l,t+n+l-1}(Y),\\
    \bar x&\in Z_{r-1}^{s,t}(Z),&
    \bar y&\in Z_\infty^{s+r,t+r-1}(Z),
  \end{aligned}
  \]
  proofs of $d_l^{h,E_{r'}}(\bar x)=x$, $d_r(\bar x)=\bar y$,
  $d_m^{g,E_{m_1+2}}(y)=\bar y$, and a proof that either the first page
  extension has no crossing or the classical differential has no crossing on
  $E_{r'}$.
  The target of the $h$-extension is an actual $E_2(\Sigma X)$
  label; the following comparison relates it to the displayed $X$ label.
\end{definition}

\begin{definition}[Actual tower suspension comparison]
  \label{def:mahowald-tower-suspension-comparison}
  \uses{def:hf2-adams-tower,def:ss-elementwise-page-presentation}
  \lean{KIP126.Classical.Adams.Suspension.TowerComparison}
  \lean{KIP126.Classical.Adams.Suspension.TowerComparison.DesuspendsClass}
  \notready
  A suspension comparison identifies every term of the actual Adams tower
  of $\Sigma X$ with the suspension of the corresponding term for $X$,
  is the identity at stage zero, and commutes with all tower maps. Its
  layer isomorphisms commute with the actual cofiber inclusions and with
  the connecting maps, including the suspension sign. Two $E_2$ labels
  are related when they have raw second-cycle representatives whose
  actual first-page maps agree after this desuspension and whose quotient
  classes are the specified labels. Existence and independence of this
  comparison have not been constructed; an arbitrary $E_2$ equivalence
  is not a substitute for these diagrams.
\end{definition}

% SOURCE: MainPaper/main.tex:1780-1917, representative choices in the proof of Theorem thm:158d451a.
\begin{lemma}[Synchronized Mahowald representatives]
  \label{lem:mahowald-synchronized-representatives}
  \uses{def:mahowald-input-data,prop:no-crossing-iff-uniform-detection,prop:classical-crossing-iff-delta-crossing,prop:page-crossing-iff-synthetic-crossing,thm:delta-classical-equivalence}
  \notready
  Under the Mahowald hypotheses, one can choose
  \[
    [\bar x]\in\pi_{t-s,t}(\nu Z/\lambda^{r'-1})
  \]
  detected by $\bar x$ such that
  $\widehat h_{r'-1}[\bar x]$ is detected by
  $\lambda^{l_1}x$ in
  $\pi_{t-s-1,t-1+e(h)}(\nu X/\lambda^{r'-1})$ and
  $\delta_{r'-1}[\bar x]$ is detected by
  $\lambda^{m_1}\bar y$ in
  $\pi_{t-s-1,t+r'-1}(\nu Z/\lambda^{m_1+1})$.
\end{lemma}

\begin{proof}
  Translate the page extension and classical differential into the two
  synthetic ESS extensions in the common quotient.  The assumed no-crossing
  of at least one extension invokes uniform detection for every representative
  chosen to realize the other.  Hence one representative realizes both
  leading terms simultaneously.
\end{proof}

% SOURCE: MainPaper/main.tex:1886-1907, the representative for y, degree-uniqueness, and the application of May's pullback lemma.
\begin{lemma}[May-compatible lift in the Mahowald diagram]
  \label{lem:mahowald-compatible-quotient-lift}
  \uses{def:mahowald-input-data,lem:mahowald-synchronized-representatives,thm:literature-route-may,prop:synthetic-ess-weightwise-transport,prop:fextension-iff-detected,thm:external-synthetic-einfty-quotient,prop:synthetic-einfty-first-quotient,prop:normalized-synthetic-triangle,def:lambda-rho-delta-triangle}
  \notready
  Under the Mahowald hypotheses, choose the synchronized representative
  $[\bar x]$ of Lemma~\ref{lem:mahowald-synchronized-representatives}.
  There are a representative

  \[
    [y]\in\pi_{t-s-1,t+n+l-1}(\nu Y/\lambda^{m_1+1})
  \]
  detected by $y$ and a class

  \[
    u\in\pi_{t-s-1,t-1+e(h)}(\nu X/\lambda^r)
  \]
  detected by $\lambda^{l_1}x$ such that

  \[
    \rho(u)=\widehat h_{r'-1}([\bar x]),
    \qquad
    \widehat f_r(u)=\lambda^{r'-1}[y].
  \]
\end{lemma}

\begin{proof}
  The fourth page-extension hypothesis and the representative criterion choose
  $[y]$ whose image under $\widehat g_{m_1+1}$ is detected by
  $\lambda^{m_1}\bar y$.  If $m_1=0$, the first-quotient collapse and its
  unique-detection conclusion make that mod-$\lambda$ homotopy class unique;
  if $m_1\ge1$, the finite-quotient $E_\infty$ formula gives uniqueness in the
  indicated top weight.  Hence in both cases
  $\widehat g_{m_1+1}([y])$ equals
  $\delta_{r'-1}([\bar x])$.  Apply May's smash-boundary theorem to the
  normalized synthetic triangle and the $\lambda$--$\rho$--$\delta$ triangle.
  Its pullback class is $u$, and its two structure maps give exactly the two
  displayed equalities.
\end{proof}

% SOURCE: MainPaper/main.tex:1900-1910, equation eq:f4848f29 and the ensuing one-way survival deduction.
\begin{lemma}[A compatible quotient lift gives source survival]
  \label{lem:mahowald-quotient-lift-survival}
  \uses{def:mahowald-input-data,lem:mahowald-synchronized-representatives,thm:external-synthetic-einfty-quotient,prop:synthetic-einfty-map-compatibility}
  \notready
  Under the Mahowald hypotheses, suppose

  \[
    u\in\pi_{t-s-1,t-1+e(h)}(\nu X/\lambda^r)
  \]
  is detected by $\lambda^{l_1}x$ and its $\rho$-image is the synchronized
  representative detected by $\lambda^{l_1}x$ modulo $\lambda^{r'-1}$.
  Then

  \[
    x\in Z_{n+m+e(h)}^{s+l,t+l-1}(X).
  \]
  No converse for a preselected representative in the smaller quotient is
  asserted.
\end{lemma}

\begin{proof}
  Apply the finite-quotient $E_\infty$ formula to the class $u$.  Its exponent
  after removing the displayed $\lambda^{l_1}$ is
  $r-l_1=n+m+e(h)$, so the existence of $u$ places $x$ in the corresponding
  cycle subgroup.  Compatibility of $\rho$ verifies that its reduction is the
  synchronized smaller-quotient representative; it is not used to manufacture
  a lift from survival alone.
\end{proof}

% SOURCE: MainPaper/main.tex:1903-1917, final lambda-division step of Theorem thm:158d451a.
\begin{lemma}[$\lambda$-division indeterminacy in the Mahowald diagram]
  \label{lem:mahowald-lambda-division}
  \uses{def:mahowald-input-data,prop:synthetic-einfty-map-compatibility,def:page-extension-coset}
  \notready
  Under the Mahowald hypotheses, if the smash diagram gives
  \[
    d_n^{\widehat f_r}(\lambda^{l_1}x)=\lambda^{r'-1}y,
  \]
  then division by $\lambda^{l_1}$ is defined in the quotient of exponent
  $r-l_1=n+m+e(h)$ and changes the target exactly by
  $B_{r'}^{s+n+l,t+n+l-1}(Y)$.  The result is
  \[
    d_n^{f,E_{n+m+1+e(h)}}(x)
       \equiv y\pmod{B_{r'}^{s+n+l,t+n+l-1}(Y)}.
  \]
\end{lemma}

\begin{proof}
  Apply the quotient-map description of multiplication by $\lambda^{l_1}$.
  Its kernel is the boundary cutoff enlarged by $l_1$; after substituting the
  definitions of $n_1,m_1,l_1,r'$, this cutoff is $B_{r'}$.  Compare with the
  complete page-extension target coset to obtain the asserted equivalence.
\end{proof}

% SOURCE: MainPaper/main.tex:1780-1917, Theorem thm:158d451a.
\begin{theorem}[Generalized Mahowald trick]
  \label{thm:generalized-mahowald}
  \uses{def:mahowald-input-data,def:mahowald-tower-suspension-comparison,prop:normalized-synthetic-triangle,def:lambda-rho-delta-triangle,lem:mahowald-synchronized-representatives,lem:mahowald-compatible-quotient-lift,lem:mahowald-quotient-lift-survival,lem:mahowald-lambda-division}
  \lean{KIP126.Kervaire.Route.Tools.GeneralizedMahowaldLaw}
  \notready
  Let
  $X\xrightarrow fY\xrightarrow gZ\xrightarrow h\Sigma X$ be distinguished
  with $e(f)+e(g)+e(h)=1$.  Put
  $r=n+m+l$, $n_1=n-e(f)\ge1$,
  $m_1=m-e(g)\ge0$, $l_1=l-e(h)\ge0$, and
  $r'=r-m_1=n_1+l_1+1$.  Suppose
  \[
  \begin{aligned}
    x&\in Z_{n_1}^{s+l,t+l-1}(X),&
    y&\in Z_{m_1+1}^{s+n+l,t+n+l-1}(Y),\\
    \bar x&\in Z_{r-1}^{s,t}(Z),&
    \bar y&\in Z_\infty^{s+r,t+r-1}(Z),
  \end{aligned}
  \]
  and assume
  \begin{enumerate}
    \item $d_l^{h,E_{r'}}(\bar x)=x$;
    \item $d_r(\bar x)=\bar y$;
    \item either (1) has no page crossing or (2) has no crossing on $E_{r'}$;
    \item $d_m^{g,E_{m_1+2}}(y)=\bar y$.
  \end{enumerate}
  Then $x\in Z_{n+m+e(h)}$ and
  \[
    d_n^{f,E_{n+m+1+e(h)}}(x)
       \equiv y\quad\text{modulo }B_{r'}.
  \]
  The linked Lean declaration is the parameterized law, using one actual
  distinguished triangle, its three normalized page families, and the
  tower suspension comparison above. A proof of this law remains to be
  supplied; defining the predicate is not a proof of the theorem.
\end{theorem}

\begin{proof}
  Smash the normalized synthetic triangle with the quotient triangle of
  exponents $r$ and $r'-1$.  The synchronization lemma chooses the common
  representative of $\bar x$.  The May-compatible-lift lemma uses condition
  (4), the degree-unique common target, and May's pullback theorem to produce a
  single class mapping to both sides of the diagram.  The one-way
  quotient-lift lemma turns that actual $\rho$ lift into the stronger survival
  statement for $x$.  Its other image is the normalized $f$ relation; divide
  by $\lambda^{l_1}$ using the final lemma to obtain the page extension with
  exactly $B_{r'}$ indeterminacy.
\end{proof}

% SOURCE: MainPaper/main.tex:1930-1977, Example exam:Mahowald; computed cofiber d_2 is external evidence.
\begin{proposition}[Cofiber regression for the Mahowald trick]
  \label{prop:mahowald-cofiber-regression}
  \uses{thm:generalized-mahowald,lem:normalized-hopf-exactness-case,def:hopf-cofiber-cell-notation,cor:classical-e2-no-crossing}
  \notready
  Given the computed differential
  $d_2(h_0h_4^2[4])=h_0p[0]$ in the Adams spectral sequence of
  $S^0/\nu_{\mathrm{Hopf}}$ and the two displayed Ext edge-map facts, the
  generalized Mahowald trick proves that $h_0h_4^2$ survives to $E_3$ and
  \[
    d_2^{\nu_{\mathrm{Hopf}},E_3}(h_0h_4^2)=h_0p.
  \]
\end{proposition}

\begin{proof}
  Use Lemma~\ref{lem:normalized-hopf-exactness-case} to instantiate the Hopf
  triangle with $(e(f),e(g),e(h))=(1,0,0)$ and its one chosen lifted triangle.
  Substitute $m=l=0$, $n=r=r'=2$, $n_1=1$, and $m_1=l_1=0$ into the theorem.
  Bottom- and top-cell Ext exactness give the two length-zero page extensions;
  the supplied cofiber differential gives condition (2), and $E_2$
  no-crossing gives condition (3).  The theorem produces the stated
  extension with zero additional indeterminacy in this degree.
\end{proof}

% SOURCE: MainPaper/main.tex:1981-2068, Proposition prop:dec738d3, first part.
\begin{proposition}[Restriction of a page extension]
  \label{prop:page-extension-restrict}
  \uses{def:page-extension,cor:fess-map-of-stable-pages,prop:lambda-rho-ess-only-d0,prop:lambda-quotient-restriction-coherence}
  \notready
  If $d_n^{f,E_r}(x)=y$ and $2\le r'<r$, reduction along
  \[
    \rho_X:\nu X/\lambda^{r-1}\to\nu X/\lambda^{r'-1},
    \qquad
    \rho_Y:\nu Y/\lambda^{r-1}\to\nu Y/\lambda^{r'-1}
  \]
  gives
  $d_n^{f,E_{r'}}(x)=y$ whenever the source and target are defined.
\end{proposition}

\begin{proof}
  The square of reduced normalized maps commutes.  Since both vertical
  $\rho$-ESS have only $d_0$, the stable-page map theorem induces a morphism of
  the two $\widehat f$ extension spectral sequences.  Naturality sends the
  length-$n$ differential to its earlier-page version.
\end{proof}

\begin{lemma}[Cokernel criterion for a specified strict solution lift]
  \label{lem:strict-solution-cokernel-obstruction}
  \uses{def:strict-ess-representative-fiber,def:filtered-chain-map}
  \lean{KIP126.Core.SpectralSequence.FilteredComplex.Solutions.liftingMap,KIP126.Core.SpectralSequence.FilteredComplex.Solutions.LiftingObstruction,KIP126.Core.SpectralSequence.FilteredComplex.Solutions.obstruction,KIP126.Core.SpectralSequence.FilteredComplex.Solutions.obstruction_eq_zero_iff}
  \leanok
  Let $f:C\to D$ be a filtered chain map, $R_C,R_D$ the actual groups
  of representative pairs, and $E_C$ the equation group of
  Definition~\ref{def:strict-ess-representative-fiber}.
  With $\rho_R:R_C\to R_D$ the induced representative map, set
  \[
    L(z)=(\Phi_C(z),\rho_R(z)),\qquad
    \operatorname{Ob}(f)=\operatorname{coker}\bigl(L:R_C\to E_C\times R_D\bigr).
  \]
  For a specified earlier solution $a$ with the images of labels $(x,y)$,
  its obstruction is the class of $((x,y,0),a)$.  That class is zero
  exactly when there is a strict solution $b$ with labels $(x,y)$ and
  $\rho_R(b)=a$.  No nonemptiness of the later fiber is assumed.
  This criterion concerns that particular representative pair $a$.
  The paper's later-relation conclusion has a different quantifier; this
  cokernel criterion does not imply a no-crossing criterion.
\end{lemma}
\begin{proof}
  \leanok
  Zero in the cokernel means that $((x,y,0),a)$ lies in the image of $L$.
  Its two coordinates are exactly the later equation and the prescribed
  restriction.  Conversely, any such solution supplies the required preimage.
\end{proof}

\begin{lemma}[Surjectivity on strict fibers and their difference groups]
  \label{lem:strict-solution-restriction-surjectivity}
  \uses{def:strict-ess-representative-fiber,def:filtered-chain-map,lem:strict-solution-cokernel-obstruction}
  \lean{KIP126.Core.SpectralSequence.FilteredComplex.Solutions.restrict_displacement,KIP126.Core.SpectralSequence.FilteredComplex.Solutions.restrict_surjective_iff_differences_surjective,KIP126.Core.SpectralSequence.FilteredComplex.Solutions.forall_obstruction_eq_zero_iff_surjective,KIP126.Core.SpectralSequence.FilteredComplex.Solutions.forall_obstruction_eq_zero_iff_differences_surjective}
  \leanok
  Let $f:C\to D$ be a filtered chain map and fix source and target
  labels for its actual strict solution fibers.  Restriction is surjective
  on these fibers if and only if the specified lifting obstruction vanishes
  for every earlier solution.  This equivalence needs no nonemptiness
  assumption.

  If a later solution $a$ is specified, surjectivity on the strict fibers is
  also equivalent to surjectivity of the actual induced homomorphism
  \[
    K_C=\ker\Phi_C\longrightarrow K_D=\ker\Phi_D.
  \]
  In particular, existence of one later solution does not establish either
  surjectivity claim.  The theorem reduces the remaining lifting obligation
  to a concrete homomorphism on the homogeneous solution groups.
\end{lemma}
\begin{proof}
  \leanok
  The first equivalence is the cokernel zero criterion applied to every
  earlier solution.  Relative to $a$ and its restriction, the coordinate
  equivalences identify a solution with its displacement from the chosen
  basepoint.  Restriction preserves displacements and translations, so
  lifting all solutions is equivalent to lifting all displacements.
\end{proof}

% SOURCE: MainPaper/main.tex:1437-1441,
% Definition def:41d51149; 1515-1536, Proposition prop:cross-f-Er; and
% 1981-2081, Proposition prop:dec738d3 and Corollary cor:dfc6043e.
% This definition includes b=0 and uses the larger ordinary boundary cutoff.
\begin{definition}[Finite nonliftable crossing candidate]
  \label{def:finite-page-extension-nonliftable-crossing}
  \uses{def:page-extension,def:page-extension-crossing,def:page-extension-coset}
  \lean{KIP126.Kervaire.Route.FiniteExtensionNonliftableCrossing}
  \leanok
  Fix a normalized family $P$, finite pages $r',r$, and integers $n,s,t$.
  A nonliftable crossing candidate consists of natural numbers $a,b$ with
  \[
    0<a\le r'-2,\qquad 0\le b\le n-a-e(f),
  \]
  classes
  \[
    x'\in Z_{r'-1-a}^{s+a,t+a}(X)
       \setminus Z_{r-1-a}^{s+a,t+a}(X),\qquad
    y'\in Z_{r'-1-n+b+e(f)}^{s+n-b,t+n-b}(Y),
  \]
  and an actual essential shorter extension in the same family,
  \[
    d_{n-a-b}^{f,E_{r'-a}}(x')=y',\qquad
    y'\notin B_{1+n-b-e(f)}^{s+n-b,t+n-b}(Y).
  \]
  The last exclusion is the larger ordinary boundary cutoff from the
  crossing definition; essentiality alone does not supply it.
  The source and target cycle conditions are those of the actual shorter
  extension witness.  The bound on $b$ also gives $a\le n-e(f)$.

  This predicate does not relate $y'$ to a chosen original target $y$.
  It enumerates all candidates in these degrees and is used only through
  their exclusion.  It is not an equivalence characterizing loss of
  essentiality, and does not describe the cokernel of a raw strict-solution
  restriction.  The project includes $b=0$ as well as $b>0$: the first
  branch of the printed proposition allows $b=0$, whereas the printed
  corollary excludes only $b>0$.
\end{definition}

\begin{definition}[Finite page-extension stretching law]
  \label{def:finite-page-extension-stretching-law}
  \uses{def:finite-page-extension-nonliftable-crossing,def:page-extension}
  \lean{KIP126.Kervaire.Route.Tools.FinitePageExtensionStretchingLaw}
  \leanok
  A normalized family $P$ satisfies the finite stretching law if, for every
  \[
    2\le r'\le r<\infty,\qquad e(f)\le n\le r'-2+e(f),
  \]
  and actual later-page labels
  \[
    x\in Z_{r-1}^{s,t}(X),\qquad
    y\in Z_{r-1-n+e(f)}^{s+n,t+n}(Y),
  \]
  the earlier relation $d_n^{f,E_{r'}}(x)=y$ and the absence of a candidate
  from Definition~\ref{def:finite-page-extension-nonliftable-crossing}
  imply $d_n^{f,E_r}(x)=y$.
  This is the precise proposition to deliver for the chosen model; defining
  it does not prove that an arbitrary normalized family satisfies it.
  Its later-page target hypothesis and its exclusion of $b=0$ candidates
  make the project statement more restrictive than printed Corollary 6.21.
  The conclusion asserts existence of a later relation, without prescribing
  the reduction of every later representative to a specified earlier pair.
\end{definition}

% SOURCE: MainPaper/main.tex:1981-2068,
% Proposition prop:dec738d3; necessary-candidate direction only, with the
% shorter-page interpretation from Proposition prop:cross-f-Er.
\begin{proposition}[Certificate for loss of essentiality on an earlier page]
  \label{prop:page-extension-loss-certificate}
  \uses{prop:page-extension-restrict,def:finite-page-extension-nonliftable-crossing,prop:page-crossing-iff-synthetic-crossing}
  \notready
  For the normalized family of the chosen synthetic model, suppose an
  essential $(f,E_r)$ extension of length $n$, with source degree $(s,t)$,
  becomes inessential on $E_{r'}$, where
  \[
    2\le r'<r<\infty,\qquad e(f)\le n\le r'-2+e(f).
  \]
  Then there is a nonliftable crossing candidate from
  Definition~\ref{def:finite-page-extension-nonliftable-crossing}.
  This is only a necessary-candidate conclusion.  An arbitrary candidate
  in those degrees is not asserted to cause loss of essentiality of the
  chosen original extension.
\end{proposition}

\begin{proof}
  The planned proof follows the printed finite descent through essential
  shorter extensions.  The first branch gives a nonliftable source with
  $b=0$; later branches may give $b>0$.  The $\lambda$ comparison converts
  the essential contribution in the fixed quotient to an actual extension
  on page $r'-a$.  This conversion must retain the larger ordinary
  boundary exclusion in
  Proposition~\ref{prop:page-crossing-iff-synthetic-crossing}.
  The model comparisons and finite descent still require canonical Lean
  proofs.  No condition identifying the arbitrary candidate target with
  the original target has been added to the candidate definition.
\end{proof}

% SOURCE: MainPaper/main.tex:1981-2081,
% Proposition prop:dec738d3 and Corollary cor:dfc6043e; project finite,
% later-typed formulation with the stronger b>=0 crossing exclusion.
\begin{lemma}[First obstruction to a later page relation]
  \label{lem:page-extension-first-obstruction}
  \uses{def:finite-page-extension-nonliftable-crossing,def:finite-page-extension-stretching-law,prop:page-extension-loss-certificate,prop:page-crossing-iff-synthetic-crossing}
  \notready
  For the normalized family of the chosen synthetic model, take the ranges
  and later-page labels of
  Definition~\ref{def:finite-page-extension-stretching-law}.  If
  $d_n^{f,E_{r'}}(x)=y$ holds but $d_n^{f,E_r}(x)=y$ does not, then there
  is a candidate from
  Definition~\ref{def:finite-page-extension-nonliftable-crossing}.
  This is a relation-level claim.  It neither asserts that every failure
  to lift a specified strict solution produces such a candidate nor
  identifies the strict-solution cokernel with the candidate list.
\end{lemma}

\begin{proof}
  The required argument is the finite descent through shorter essential
  extensions in the actual $\widehat f$-ESS, using the $\rho$ and
  $\lambda$ comparison maps.  The comparison with a shorter page must
  preserve the larger ordinary boundary exclusion, as in
  Proposition~\ref{prop:page-crossing-iff-synthetic-crossing}.
  The canonical proof of this relation-level implication remains to be
  supplied.  Proposition 6.20 discusses loss of essentiality of a known
  later relation, and the printed Corollary 6.21 states existence of a
  later relation; their passage must be checked with the stated ranges.
  The generic cokernel criterion does not prove this step, and no
  surjectivity assertion about raw strict representatives is used here.
\end{proof}

% SOURCE: MainPaper/main.tex:2074-2081,
% Corollary cor:dfc6043e.  The infinite case additionally requires the
% project coherent-tower and untruncated comparison obligations below.
\begin{corollary}[Page-extension stretching]
  \label{cor:page-extension-stretch}
  \uses{def:finite-page-extension-stretching-law,lem:page-extension-first-obstruction,def:synthetic-ess-solution-tower,lem:synthetic-ess-coherent-tower-limit}
  \notready
  The normalized family of the chosen synthetic model satisfies the finite
  stretching law of
  Definition~\ref{def:finite-page-extension-stretching-law}.
  For permanent labels $x,y$, an earlier extension in the same range
  therefore yields the corresponding relation on every later finite page
  if all the nonliftable crossing candidates are excluded.

  For the untruncated conclusion $d_n^{f,E_\infty}(x)=y$, the finite
  strict solutions must additionally admit a coherent tower and the
  actual untruncated comparison of
  Lemma~\ref{lem:synthetic-ess-coherent-tower-limit} must be established
  for the model.  Explicit surjectivity of the actual adjacent restriction
  maps is one sufficient condition for constructing such a tower from an
  initial solution.  Alternatively, finiteness of every actual strict
  fiber, together with their nonemptiness, gives some coherent tower by
  Proposition~\ref{prop:synthetic-ess-coherent-tower-of-finite}.
  Exclusion of the finite candidates is not asserted to imply either
  surjectivity or finiteness of those representative sets.

  The two root predicates now fix the finite candidate list and the
  finite law, but the law's proof and its model comparison witnesses
  remain outstanding.  The untruncated comparison is a further proof
  obligation.  The retired free-operations statement expressed none of
  these conclusions.
\end{corollary}

\begin{proof}
  The finite law is the contrapositive of the first-obstruction lemma,
  whose model proof remains open.  The resulting finite relations give
  nonempty strict solution fibers.  To pass to $E_\infty$, construct a
  coherent tower using a separately justified compatibility argument and
  then prove that its limit gives an untruncated solution for $\widehat f$.
  Nonemptiness of each finite fiber alone completes neither step.
\end{proof}

% SOURCE: MainPaper/main.tex:2084-2099,
% Example exam:stretchext.  The printed example does not assert
% surjectivity of restrictions on all raw strict representative pairs.
\begin{proposition}[Stretching the Hopf-map extension to $E_\infty$]
  \label{prop:page-stretch-regression}
  \uses{prop:hopf-e3-page-extension,evidence:hopf-crossing-obstruction-exclusion,cor:page-extension-stretch,def:synthetic-ess-solution-tower,lem:synthetic-ess-coherent-tower-limit}
  \notready
  The $E_3$ extension
  $d_2^{\nu_{\mathrm{Hopf}},E_3}(h_0h_4^2)=h_0p$ stretches to the
  essential extension
  $d_2^{\nu_{\mathrm{Hopf}},E_\infty}(h_0h_4^2)=h_0p$.
  Its finite crossing exclusion supplies no assertion that every specified
  strict representative pair lifts between adjacent quotients.
\end{proposition}

\begin{proof}
  The independent finite Hopf certificate excludes the degree-allowed
  crossing candidates.  It also supplies that $h_0h_4^2$ and $h_0p$ are
  permanent representatives, giving both later-page label types.
  Once the finite stretching law is proved, these facts give an extension
  on every finite page.  To obtain the displayed untruncated conclusion,
  a coherent tower must still be constructed and compared to an actual
  untruncated solution, or an independent untruncated lifting argument
  must be supplied.  The finite candidate exclusion alone does not fill
  this gap.  Finally, the certificate's nonzero complete target-coset
  clause makes the resulting length-two relation essential, once its
  existence is established.
\end{proof}

% SOURCE: MainPaper/main.tex:1398-1405 and 929-961; project packaging of the
% zero-homology rotation needed by the normalized-cofiber interface.
\begin{lemma}[Normalized Hopf cofiber exactness case]
  \label{lem:normalized-hopf-exactness-case}
  \uses{evidence:normalized-hopf-detection,def:normalized-synthetic-exactness-case,prop:positive-adams-filtration-homology-zero,def:cofiber-triangle,def:hf2-homology,thm:external-synthetic-triangle-lift}
  \notready
  The classical Hopf cofiber triangle
  \[
    S^3\xrightarrow{\nu_{\mathrm{Hopf}}}S^0\xrightarrow g
    S^0/\nu_{\mathrm{Hopf}}\xrightarrow hS^4
  \]
  carries the zero constructor of
  Definition~\ref{def:normalized-synthetic-exactness-case}.  Explicitly,
  \[
    H_*\nu_{\mathrm{Hopf}}=0,
    \qquad
    (e(\nu_{\mathrm{Hopf}}),e(g),e(h))=(1,0,0),
  \]
  and its long exact homology sequence supplies the rotated short exact
  sequence.  The constructor's lifted-triangle field is exactly the witness
  retained by Proposition~\ref{evidence:normalized-hopf-detection}, and its
  normalized Hopf component is therefore equal to $[h_2]$.
\end{lemma}

\begin{proof}
  The filtration-one detection by $h_2$ gives
  $e(\nu_{\mathrm{Hopf}})=1$, while degree considerations make the Hopf map
  zero on $H\mathbb F_2$-homology.  In the cofiber long exact sequence, $g$
  is the nonzero bottom-cell inclusion on homology and $h$ is the nonzero
  top-cell quotient.  The contrapositive of
  Proposition~\ref{prop:positive-adams-filtration-homology-zero} therefore
  gives Adams filtration zero and $e(g)=e(h)=0$.  The remaining
  long-exact-sequence terms split by degree into
  the required short exact sequence.  Package these facts and the displayed
  $e$-pattern around the precise lifted-triangle witness retained by the
  normalized-detection evidence; do not invoke the nonunique lift theorem a
  second time.  The evidence's component equality becomes the zero tag's
  equality $\widehat{\nu_{\mathrm{Hopf}}}=[h_2]$.
\end{proof}
\begin{theorem}[Local sphere specialization of the general Moss input]
  \label{thm:main-local-moss-specialization}
  \uses{thm:literature-moss,def:adams-hj,def:adams-detection,def:toda-bracket}
  \lean{KIP126.Main.Solution.Literature.route_moss}
  \notready
  The general sphere Moss theorem specializes to the local $E_3$ implication
  for $\langle h_5^2,h_0,B\rangle$. The actual mapping tower,
  composition, Milnor labels and
  source convergence must be compared. All defining-system, detection,
  null-product, noncrossing and residual-injectivity conditions remain
  explicit. This is an internal theorem with an unfinished proof.
\end{theorem}

\begin{theorem}[Naturality of the fixed synthetic weight comparison]
  \label{thm:main-weight-shift-naturality}
  \uses{def:synthetic-adams-ss,def:synthetic-tridegree,def:hf2-adams-tower}
  \lean{KIP126.Synthetic.SpectralSequence.weightShiftObjectIso,KIP126.Synthetic.SpectralSequence.canonicalWeightShiftAmbient,KIP126.Synthetic.SpectralSequence.CanonicalWeightShift,KIP126.Synthetic.SpectralSequence.canonicalWeightShift_exists,KIP126.Def.standardEInftyWeightShift,KIP126.Main.Solution.Literature.eInfty_shift_natural}
  \notready
  The actual object weight isomorphism and the fixed tower presentation
  prescribe the map on permanent cycles. The existence theorem
  must descend this map and its inverse to $E_\infty$. This fixed comparison
  is natural on all synthetic objects. Neither its
  construction nor full-category naturality is an independently selected
  BHS literature input; both property proofs remain unfinished.
\end{theorem}

\begin{theorem}[BR21 in the chosen tmf coordinates]
  \label{thm:main-tmf-br21-coordinate-differential}
  \uses{thm:literature-br21,def:tmf-csv-coordinate-binding,def:hf2-adams-tower,thm:computation-bindings-tmf-v2Sixteen,thm:computation-bindings-tmf-betaGFour,thm:computation-tmfMultiplicative}
  \lean{KIP126.Main.Solution.Computation.tmf_br21_d3}
  \notready
  Transport the BR21 differential
  $d_3(w_2^2)=\beta g^4$ through its actual detector comparison and use
  its computation coordinates and the proved relation
  $\beta g^4=\beta^5g$. The resulting equation is
  $d_3(v_2^{16})=\beta^5g$ in bidegrees
  $(16,112)\to(19,114)$. The transport proof is implemented and uses the
  existing unfinished generic Adams tower naturality theorem. It is not
  an additional computation input or a completed proof claim.
\end{theorem}
